%% =====================================================================
%%  Capítulo 1 -- Introducción: qué es el análisis de sistemas de control
%%  y de dónde viene.
%%  Migrado al formato nuevo (noviembre de 2026) desde Introduccion.tex.
%%  Cambios de fondo respecto a la versión anterior (ver MIGRACION.md):
%%    la historia se reorganizó por épocas, se verificaron fechas y nombres
%%    contra fuentes primarias o reseñas especializadas y se corrigieron
%%    varios datos (Watt no patentó el regulador; Sperry voló en 1914;
%%    Bode publicó en 1940--45; IFAC se fundó en 1957; RIAS en 1955...);
%%    el tanque en lazo abierto ahora se compara con el lazo cerrado,
%%    con números verificados; "malla abierta" y "lazo abierto" se
%%    distinguen de "control en lazo abierto".
%% =====================================================================

\capitulo[cap:in]{Introducción: el análisis de sistemas de control}
  {Es que con los ojos del presente ningún pasado es tolerable\ldots{} Nada hay más ingrato e implacable que la eternidad; nada se purga y se corrige con más violencia que la historia.}
  {Juan Esteban Constaín}

\begin{objetivos}[partes/01-introduccion.html]
  \item Explicar qué se entiende por \emph{análisis}, \emph{sistema} y \emph{control}, y qué hace la retroalimentación.
  \item Distinguir el control en lazo abierto del control en lazo cerrado y medir, con un ejemplo, qué gana y qué cuesta retroalimentar.
  \item Plantear un problema de control: objetivos, variables manipuladas y perturbaciones, mediciones, restricciones y estructura del controlador.
  \item Ubicar en el tiempo las ideas del curso, desde los reguladores de flotador de la Antigüedad hasta el control digital, y saber dónde leer más.
  \item Saber cómo están organizados el curso, las notas, los códigos y la página web.
\end{objetivos}

%% ---------------------------------------------------------------------
\section{Análisis, sistema y control}
\label{sec:in-palabras}
%% ---------------------------------------------------------------------

El nombre del curso tiene tres palabras y cada una carga con un significado preciso. Vale la pena detenerse en ellas antes de escribir la primera ecuación.

\paragraph{Análisis}
El diccionario de la Real Academia da varias acepciones de \emph{análisis}: la distinción y separación de las partes de algo para conocer su composición, el estudio detallado de algo y, la que nos interesa, la \emph{parte de las matemáticas basada en los conceptos de límite, convergencia y continuidad}, que da origen al cálculo diferencial e integral y a la teoría de funciones. En el curso las dos ideas van juntas: \emph{separamos} un sistema en sus partes (planta, actuador, sensor, controlador) y \emph{estudiamos} cada una con herramientas del análisis matemático: ecuaciones diferenciales, transformada de Laplace, variable compleja.

\paragraph{Sistema}
Un sistema es un conjunto de cosas (personas, células, moléculas, tuberías, transistores) interconectadas de tal forma que producen un patrón de comportamiento a medida que pasa el tiempo; puede ser activado, restringido o impulsado por fuerzas externas, pero su respuesta depende de sus propias características~\citep{MeWr08}. Más corto: una combinación de componentes que se percibe como un todo~\citep{HaCh08}. Lo que hace interesantes a los sistemas que estudiaremos es que son \emph{dinámicos}: lo que hacen ahora depende de lo que les pasó antes. En el capítulo~\ref{cap:po} precisaremos estas ideas con las nociones de señal y de sistema dinámico.

\paragraph{Control}
La Real Academia también da varias acepciones de \emph{control} (del francés \emph{contrôle}): comprobación, dominio, el examen parcial con el que se comprueba la marcha de los alumnos\ldots{} y la que nos interesa: \emph{regulación, manual o automática, sobre un sistema}.

\begin{definicion}[label=def:in-sistemacontrol]{Sistema de control}
Un \clave{sistema de control} es un conjunto de dispositivos que administra, ordena, dirige o regula el comportamiento de otro sistema para obtener un resultado deseado~\citep{WikiControl}. Está formado por subsistemas (el proceso o planta, los actuadores, los sensores, el controlador) ensamblados para que una salida siga una referencia o rechace perturbaciones, con un desempeño especificado.
\end{definicion}

El \emph{análisis de sistemas de control}, entonces, estudia cómo evolucionan en el tiempo esos elementos interconectados y qué tan bien cumplen su propósito: si el lazo es estable, qué tan rápido y con cuánto sobrepaso responde, cuánto error deja, qué tanto le afectan las perturbaciones y los errores del modelo.

\subsection{La idea central: retroalimentar}

\begin{definicion}[label=def:in-retro]{Retroalimentación}
\clave{Retroalimentar} es medir la variable que se quiere controlar, compararla con el valor deseado y usar la diferencia (el \emph{error}) para decidir la acción. Un sistema así se llama \clave{retroalimentado} o \emph{de lazo cerrado}.
\end{definicion}

La retroalimentación es el concepto fundamental del curso y su gran virtud es la \clave{robustez}: como la acción depende de lo que realmente pasa, el sistema puede cumplir sus requisitos aunque el modelo sea imperfecto, aparezcan perturbaciones o las mediciones lleguen con ruido. Su peligro es la \alerta{sobrecorrección}: si se reacciona demasiado fuerte o demasiado tarde, el lazo oscila y puede volverse inestable. Casi todo el curso trata de cómo aprovechar la virtud sin caer en el peligro.

Richard Murray resumió el área en su conferencia Bode de 2016 con una ecuación de cuatro términos, \emph{control = dynamics, uncertainty, feedforward, feedback}~\citep{Mur16}, cuyas iniciales (DUFF) recuerdan la cerveza de \emph{Los Simpson}. Leída en orden: (i) hay unas \clave{dinámicas} que queremos modificar; (ii) las conocemos con \clave{incertidumbre}; (iii) con la \clave{prealimentación} (\emph{feedforward}) actuamos con lo que sabemos de antemano: la referencia, o una perturbación que podemos medir; y (iv) con la \clave{retroalimentación} (\emph{feedback}) corregimos lo que no sabíamos.

\begin{ideaclave}
\textbf{Señales y sistemas.}
La retroalimentación se entiende con dos conceptos entrelazados~\citep{AlMa10}: las \emph{señales}, porque es midiéndolas como se decide qué hacer, y los \emph{sistemas}, porque la retroalimentación existe para cambiar su comportamiento. Cuando hablamos de un sistema hablamos de algo que produce o modifica una señal; cuando hablamos de una señal, de algo que algún sistema produjo.
\end{ideaclave}

\begin{pausa}
Cuando usted conduce un carro por una carretera con curvas, ¿quién hace de sensor, de controlador y de actuador? ¿Qué es la referencia y qué la perturbación?
\respuesta{Los ojos (y el oído interno) son el sensor; el cerebro compara la posición del carro con el centro del carril (la referencia) y decide: es el controlador; los brazos y el timón son el actuador. El viento lateral, los baches y la inclinación de la vía son perturbaciones. Si usted cierra los ojos, el lazo se abre.}
\end{pausa}

%% ---------------------------------------------------------------------
\section{Lazo abierto y lazo cerrado}
\label{sec:in-lazos}
%% ---------------------------------------------------------------------

Hay dos maneras de mandar sobre un sistema (\cref{fig:in-lazos}). En el \clave{control en lazo abierto} el controlador decide la acción sin mirar la salida: un actuador mueve el proceso según una calibración hecha de antemano. En el \clave{control en lazo cerrado} un sensor mide la salida, el comparador calcula el error $e=r-y_m$ y el controlador actúa sobre ese error.

\begin{figure}[ht]
\centering
\includegraphics[width=\linewidth]{cap01/lazos_ab}
\caption{(a) Control en lazo abierto: nadie mira la salida. (b) Control en lazo cerrado: el sensor cierra el lazo. Notación del curso: $r$ referencia (\emph{set point}), $e$ error, $u$ señal de control, $d$ perturbación, $y$ variable controlada, $y_m$ medición, $n$ ruido de medición.}
\label{fig:in-lazos}
\end{figure}

Una tostadora, un semáforo de tiempos fijos o una lavadora que sigue su programa trabajan en lazo abierto: son baratos y no necesitan sensor, pero dependen de que el mundo se comporte como se supuso al calibrarlos. Un termostato, el control de crucero de un carro o el regulador de una máquina de vapor trabajan en lazo cerrado. El ejemplo siguiente pone números a la diferencia.

\begin{ejemplo}[label=ej:in-tanque]{Llenar un tanque sin mirarlo y mirándolo}
Un tanque de área $A=0.5$\,m$^2$ recibe un caudal $u(t)$ (m$^3$/min) que fijamos con una válvula, puede recibir lluvia $d(t)$ y descarga por una válvula de resistencia $R=2$\,min/m$^2$. Alrededor del punto de operación, el nivel $h$ (m) cumple
\begin{equation}
  A\,\frac{\dif h}{\dif t}=u+d-\frac{h}{R}
  \quad\Longrightarrow\quad
  \tau\frac{\dif h}{\dif t}+h=K\,(u+d),\qquad K=R=2,\quad \tau=AR=1\text{ min}.
  \label{eq:in-tanque}
\end{equation}
Queremos $h=h_r=1$\,m. Compare dos estrategias: (i) lazo abierto calibrado, $u=h_r/K$; (ii) lazo cerrado, $u=h_r/K+K_c\,(h_r-h)$ con $K_c=5$. Evalúe el nivel final en tres casos: sin perturbación, con lluvia $d=0.2$\,m$^3$/min y con la válvula de salida parcialmente obstruida ($R$ sube un $20\%$, a $2.4$).
\solucion En estado estacionario $\dif h/\dif t=0$ y \eqref{eq:in-tanque} da $h=R\,(u+d)$.

\emph{Lazo abierto.} Con $u=h_r/2=0.5$: sin perturbación $h=2(0.5)=1$, perfecto; con lluvia $h=2(0.5+0.2)=1.4$\,m; con la válvula obstruida $h=2.4(0.5)=1.2$\,m. El lazo abierto es exacto \emph{solo} si el modelo es exacto y no hay perturbaciones.

\emph{Lazo cerrado.} Ahora $u$ depende de $h$: $h=R\bigl(0.5+K_c(1-h)+d\bigr)$, de donde
\begin{equation}
  h=\frac{R\,(0.5+K_c+d)}{1+RK_c}.
  \label{eq:in-tanquecerrado}
\end{equation}
Sin perturbación $h=1$; con lluvia $h=11.4/11=1.036$\,m; con la válvula obstruida $h=2.4(5.5)/13=1.015$\,m. La lluvia movía el nivel $0.4$\,m en lazo abierto y ahora lo mueve $0.4/11=0.036$\,m: \emph{once veces menos}. El $11$ es $1+KK_c$, y al producto $KK_c=10$ lo llamaremos \clave{ganancia de lazo}. La \cref{fig:in-tanque} muestra las respuestas completas.
\end{ejemplo}

\begin{figure}[ht]
\centering
\includegraphics[width=\linewidth]{cap01/f_tanque_lazos}
\caption{Ejemplo~\ref{ej:in-tanque}. El lazo abierto (coral) llega exacto mientras el mundo coincide con la calibración; la retroalimentación (azul) reduce el efecto de la lluvia y de la obstrucción por un factor $1+KK_c$ y, de paso, responde más rápido.}
\label{fig:in-tanque}
\end{figure}

\codigorepo{cap01_introduccion/C1_Ej1_TanqueLazos.m}{cerrado}{ejemplo \ref{ej:in-tanque}}

Hay un segundo regalo, menos evidente: el lazo cerrado es más \emph{rápido}. Si se reemplaza la ley de control en \eqref{eq:in-tanque} se obtiene $\tau\dot h+(1+KK_c)\,h=\ldots$, una constante de tiempo de $\tau/(1+KK_c)=1/11$\,min$\approx5.5$\,s en lugar de $1$\,min. Nada es gratis, sin embargo: en el instante del escalón la válvula pasa de $0.5$ a $0.5+K_c=5.5$\,m$^3$/min (\cref{fig:in-esfuerzo}). Si la válvula no da más de, digamos, $2$\,m$^3$/min, se \emph{satura} y la respuesta real será más lenta que la calculada.

\begin{figure}[ht]
\centering
\includegraphics[width=0.7\linewidth]{cap01/f_tanque_esfuerzo}
\caption{El precio de la retroalimentación del ejemplo~\ref{ej:in-tanque} con lluvia: el caudal pedido salta a $5.5$\,m$^3$/min y luego se reduce para compensar la lluvia.}
\label{fig:in-esfuerzo}
\end{figure}

\begin{ideaclave}
\textbf{Lo que compra la retroalimentación.}
Con una ganancia de lazo $KK_c$: las perturbaciones y los errores del modelo se reducen aproximadamente por $1+KK_c$, y la respuesta se acelera en la misma proporción. El precio: más esfuerzo de control, un sensor (con su ruido) y el riesgo de inestabilidad cuando la planta tiene más dinámica que la de este tanque. Con un controlador proporcional queda un error pequeño pero no nulo; anularlo es trabajo de la acción integral (capítulos~6 a~8).
\end{ideaclave}

\begin{modelo}[TanqueLazos.slx]
\texttt{crearModelo\_TanqueLazos.m} construye el modelo de la versión anterior de estas notas, ahora con las dos ramas: el mismo tanque en lazo abierto (bloque \texttt{Calibracion}, que usted ajusta a mano) y en lazo cerrado (bloque \texttt{Kc}). La lluvia se activa con \texttt{Lluvia} $=1$ y la obstrucción cambiando \texttt{R} en el espacio de trabajo. Primero ajuste \texttt{Calibracion} para que el nivel llegue a $1$ sin lluvia; luego active la lluvia e intente, \emph{a mano}, mantener el nivel. ¿Pudo? Compare con la rama retroalimentada.
\end{modelo}

\begin{pausa}
En el ejemplo~\ref{ej:in-tanque}, ¿qué $K_c$ hace falta para que la lluvia mueva el nivel menos de $1$\,cm? ¿Cuánto vale entonces el caudal en el instante del escalón?
\respuesta{La desviación es $Kd/(1+KK_c)=0.4/(1+2K_c)<0.01$, es decir $K_c>19.5$. Con $K_c=20$ la desviación es $0.0098$\,m y el caudal inicial es $0.5+20=20.5$\,m$^3$/min: cuarenta veces el de lazo abierto. Pedir mucha precisión con un controlador proporcional sale caro.}
\end{pausa}

\subsection{Tres nombres parecidos}
\label{sec:in-nombres}

En español usamos ``lazo'' y ``malla'' para cosas distintas, y conviene fijarlas desde el principio porque los libros en inglés usan \emph{open loop} para las tres.

\begin{definicion}[label=def:in-nombres]{Control en lazo abierto, malla abierta y lazo abierto}
\begin{itemize}[leftmargin=1.2em]
  \item \clave{Control en lazo abierto} (\emph{open-loop control}): una \emph{estrategia} sin retroalimentación, como la de la \cref{fig:in-lazos}(a).
  \item \clave{Malla abierta} $\Gma(s)$: en un lazo cerrado, el producto de todo lo que recorre una señal al dar la vuelta completa, sin el signo del comparador. Si el camino directo es $G(s)$ y la retroalimentación $\Hs(s)$, $\Gma(s)=G(s)\Hs(s)$. Es la que se usa para estudiar estabilidad y la que se dibuja en el diagrama de Bode (en inglés, \emph{loop transfer function}, $L$).
  \item \clave{Lazo abierto} $\Gla(s)$: la función que relaciona el error con la salida cuando el lazo se reescribe con retroalimentación unitaria. Se usa para el tipo del sistema y los errores en estado estacionario.
\end{itemize}
\end{definicion}

Con $\Hs(s)=1$ la malla abierta y el lazo abierto coinciden, y por eso muchos textos los confunden; en el capítulo~\ref{cap:rf} (\cref{sec:rf-lenguaje}) se ve por qué no son lo mismo cuando el sensor tiene dinámica. La palabra ``abierto'' en los dos últimos nombres no quiere decir que el sistema funcione sin retroalimentación: describe qué parte del lazo cerrado se está mirando.

%% ---------------------------------------------------------------------
\section{Desempeño, análisis y diseño}
\label{sec:in-desempeno}
%% ---------------------------------------------------------------------

¿Cómo se juzga si un sistema de control es bueno? Tomemos un ascensor~\citep{Nis11}. Oprimir el botón del cuarto piso es la \emph{entrada}; la posición de la cabina es la \emph{salida}. El ascensor arranca, acelera, frena y se detiene. Nos importan tres cosas (\cref{fig:in-ascensor}):

\begin{itemize}
  \item La \clave{respuesta transitoria}: cómo llega. Si es muy rápida, el pasajero sufre aceleraciones incómodas o un sobrepaso (la cabina se pasa del piso y regresa); si es muy lenta, se agota su paciencia.
  \item El \clave{error en estado estacionario}: dónde se queda. Si la cabina se detiene unos centímetros por debajo del piso, alguien se va a tropezar.
  \item La \clave{estabilidad}: que se quede. Un ascensor que oscila alrededor del piso sin detenerse no sirve, por bueno que sea lo demás.
\end{itemize}

\begin{figure}[ht]
\centering
\includegraphics[width=0.85\linewidth]{cap01/f_ascensor}
\caption{Respuesta de un ascensor (modelo de segundo orden con valores exagerados para que se vean): sobrepaso de $9.5\%$, tiempo de establecimiento de $7.4$\,s y un error final de $0.1$ piso.}
\label{fig:in-ascensor}
\end{figure}

\codigorepo{cap01_introduccion/C1_Ej2_Ascensor.m}{ascensor}{respuesta del ascensor de la \cref{fig:in-ascensor}}

Estas tres cosas tienen nombre y número, y aprender a calcularlos es buena parte del curso: la estabilidad en el capítulo~5, el transitorio de primer y segundo orden en los capítulos~2 y~3 y el error en el capítulo~6.\nota{El ejemplo del ascensor tiene un video en el MOOC \emph{Introducción al análisis de sistemas de control con MATLAB}, en Coursera: \url{https://www.coursera.org/learn/introduccion-al-analisis-de-sistemas-de-control-con-matlab}.}

\begin{definicion}[label=def:in-diseno]{Análisis y diseño}
El \clave{análisis} determina el desempeño de un sistema dado: su estabilidad, su respuesta transitoria y su error. El \clave{diseño} es el medio por el cual ese desempeño se logra o se modifica: cambiando parámetros del sistema o agregando un subsistema nuevo (un controlador, un compensador) hasta cumplir las \clave{especificaciones}.
\end{definicion}

Análisis y diseño son dos direcciones del mismo camino: el análisis va del sistema al desempeño; el diseño, del desempeño deseado al sistema. Por eso el curso se llama \emph{análisis} pero termina diseñando controladores: no se puede diseñar lo que no se sabe analizar.

\begin{pausa}
El fabricante quiere eliminar el sobrepaso del ascensor de la \cref{fig:in-ascensor} sin cambiar nada más. ¿Qué cree que le pasará al tiempo de llegada?
\respuesta{Aumentará. Para no pasarse, el ascensor debe frenar antes y llegar ``despacio'' al piso; con la misma fuerza disponible tarda más. Transitorio rápido y sin sobrepaso es un compromiso, que en el capítulo~3 se cuantifica con el factor de amortiguamiento $\zeta$.}
\end{pausa}

%% ---------------------------------------------------------------------
\section{Cómo se plantea un problema de control}
\label{sec:in-problema}
%% ---------------------------------------------------------------------

\subsection{Entradas, salidas y mediciones}

Un sistema se separa de su entorno con una frontera imaginaria; lo que cruza la frontera hacia adentro son \emph{entradas}, lo que sale e influye sobre el entorno son \emph{salidas}~\citep{Bol15}. En los diagramas de bloques la frontera es la caja y las señales son las flechas. Pero para controlar hace falta distinguir más: no todas las entradas se pueden mover ni todas las salidas se pueden medir (\cref{fig:in-entradas}).

\begin{figure}[ht]
\centering
\includegraphics[width=0.8\linewidth]{cap01/proceso_entradas}
\caption{Un proceso con sus entradas y salidas clasificadas (adaptado de~\citep{Beq03}).}
\label{fig:in-entradas}
\end{figure}

\begin{definicion}[label=def:in-variables]{Variable manipulada y perturbación}
Una \clave{variable manipulada} es una entrada que el sistema de control (o un operador) puede ajustar. Una \clave{perturbación} es una entrada que afecta la salida pero que el sistema de control no puede ajustar; a veces se puede medir y a veces no. La salida que se quiere llevar a un valor deseado es la \clave{variable controlada}.
\end{definicion}

Lo que el controlador ``ve'' de la salida llega a través de un sistema de instrumentación con tres partes: el \emph{sensor}, en contacto con la variable física; el \emph{procesamiento de señal}, que convierte la salida del sensor en algo que se pueda transmitir (en la industria, típicamente una corriente de $4$ a $20$\,mA); y la \emph{presentación} o almacenamiento del dato. Al conjunto sensor--acondicionador se le llama a menudo \clave{transductor}: convierte una variable física en otra que se puede usar, casi siempre eléctrica. La medición nunca es exactamente la variable: hay errores de calibración, ruido y dinámica del propio sensor. Por eso en la \cref{fig:in-lazos} distinguimos $y$ de $y_m$.

\subsection{Los elementos de una estrategia de control}

\begin{receta}[Cómo plantear un problema de control~\citep{Beq03}]
\begin{enumerate}
  \item \textbf{Objetivos.} ¿Qué se quiere lograr? Mantener una variable en un valor, seguir una trayectoria, minimizar un costo.
  \item \textbf{Entradas.} Clasificarlas en manipuladas y perturbaciones. Pueden variar de forma continua o cambiar solo en ciertos instantes.
  \item \textbf{Salidas.} Clasificarlas en medidas y no medidas.
  \item \textbf{Restricciones.} \emph{Duras}, que no se pueden violar (un caudal entre un mínimo y un máximo, una válvula entre $0$ y $100\%$), y \emph{blandas}, que conviene respetar pero cuya violación no es peligrosa (la composición de un producto).
  \item \textbf{Operación.} \emph{Continua}, con periodos largos en condiciones parecidas, o \emph{por lotes} (\emph{batch}), con periodos cortos en los que las condiciones cambian mucho, como la fermentación de la cerveza.
  \item \textbf{Seguridad, ambiente y economía.} Qué no puede pasar nunca, qué se puede emitir, cuánto cuesta cada unidad de energía o de producto.
  \item \textbf{Estructura del controlador.} \emph{Retroalimentación}: se mide la variable controlada, se compara con la referencia y se ajusta la manipulada. \emph{Prealimentación}: se mide la perturbación y se ajusta la manipulada \emph{antes} de que la salida se afecte.
\end{enumerate}
\end{receta}

Escoger qué mover y qué medir no es trivial. Una buena \textbf{entrada de control} está \emph{disponible} (influye directa e inmediatamente sobre la salida), \emph{responde rápido}, permite ajustes \emph{finos} y tiene un efecto de \emph{signo conocido} y constante: si en un punto de operación abrir la válvula sube el nivel y en otro lo baja, el controlador no sabe para dónde ir. Una buena \textbf{salida de control} se puede medir de forma precisa, rápida y confiable, es \emph{relevante} (una buena medida de lo que realmente nos importa), responde rápido a lo que pasa en el proceso y es \emph{suave}, es decir, no necesita filtros fuertes.

\begin{definicion}[label=def:in-ganancia]{Ganancia del proceso}
La \clave{ganancia del proceso} mide la sensibilidad de la salida a la entrada: cuánto cambia la salida en estado estacionario por unidad de cambio en la entrada. Si un aumento de $0.5$\,kW en la potencia de un calentador sube la temperatura de un fluido $10$\,°C, la ganancia es $10/0.5=20$\,°C/kW.
\end{definicion}

La ganancia tiene unidades y signo, y ambos importan: en el ejemplo~\ref{ej:in-tanque} era $K=2$\,m por m$^3$/min. Una vez escogida la estructura, falta el \emph{algoritmo} del controlador y sus parámetros. En este curso los escogeremos \emph{a partir de un modelo} del proceso: por eso los capítulos~2 a~4 están dedicados a construir modelos.

\subsection{Ejemplo: darse una ducha}

Darse una ducha es un problema de control multivariable que resolvemos a diario (\cref{fig:in-ducha}).

\begin{figure}[ht]
\centering
\includegraphics[width=\linewidth]{cap01/ducha}
\caption{La ducha como problema de control: dos válvulas, dos variables que importan (temperatura y caudal) y un sensor muy bueno, la piel.}
\label{fig:in-ducha}
\end{figure}

\begin{itemize}
  \item \textbf{Objetivos:} estar limpio, refrescarse y sentirse bien; tomamos el último: una temperatura y un caudal agradables.
  \item \textbf{Manipuladas:} la apertura de las válvulas de agua caliente y fría (y, si se quiere, la posición del cuerpo).
  \item \textbf{Perturbaciones:} caídas de presión, cambios en la temperatura del calentador.
  \item \textbf{Salidas medidas:} temperatura y caudal del agua cuando toca la piel.
  \item \textbf{Restricciones:} las válvulas abren entre un mínimo y un máximo; por lo tanto hay temperaturas mínima y máxima alcanzables.
  \item \textbf{Estructura:} es \emph{multivariable}: cada válvula mueve la temperatura \emph{y} el caudal, así que hay que coordinarlas. Retroalimentamos cuando la piel siente frío y abrimos la caliente; prealimentamos cuando oímos que alguien descarga el sanitario y, antes de quemarnos, cerramos un poco la caliente.
\end{itemize}

\subsection{Ejemplo: un tanque de almacenamiento}
\label{sec:in-surge}

En una planta, los tanques de almacenamiento (\emph{surge tanks}) amortiguan las diferencias entre lo que produce una unidad y lo que consume la siguiente~\citep{Beq03}. Entra un caudal $F_1$ desde el proceso~1 y sale un caudal $F_2$ hacia el proceso~2. El \textbf{objetivo} es mantener el nivel $h$ dentro de unos límites alrededor de un valor deseado $h_{sp}$ (\emph{set point}); las \textbf{entradas} son $F_1$ y $F_2$; la \textbf{salida} es el nivel, que suponemos medible; la \textbf{operación} es continua; y la \textbf{seguridad} pide no rebosar ni vaciar el tanque y que el fluido no dañe tuberías ni paredes. Cuál entrada es manipulada y cuál perturbación \emph{depende de quién manda} (\cref{fig:in-escenarios}).

\begin{figure}[ht]
\begin{anchocompleto}
\centering
\includegraphics[width=\linewidth]{cap01/tanque_escenarios}
\end{anchocompleto}
\caption{Tres estructuras para el mismo tanque, en símbolos de la norma ISA-5.1: LT transmisor de nivel, LC controlador de nivel, FT transmisor de flujo, FY cálculo o relé; en línea discontinua, las señales.}
\label{fig:in-escenarios}
\end{figure}

\paragraph{Escenario 1: el proceso 2 fija $F_2$}
Si el proceso~2 es, por ejemplo, una caldera que toma el agua que necesita según la demanda de vapor, $F_2$ es una perturbación para el tanque y la variable manipulada es $F_1$. El transmisor LT envía el nivel al controlador LC, que lo compara con $h_{sp}$ y manda una señal de presión $P_v$ a la válvula de entrada. La \cref{fig:in-esquemas}(a) muestra el mismo lazo como diagrama de bloques. Note la diferencia: el diagrama de instrumentación dibuja \emph{equipos}; el diagrama de bloques dibuja \emph{flujo de información}, y cada bloque es un sistema dinámico.

\paragraph{Escenario 2: el proceso 1 fija $F_1$}
Si el proceso~1 produce un compuesto que el proceso~2 consume, $F_1$ es la perturbación y se manipula $F_2$ con la válvula de salida. El diagrama de bloques es el mismo intercambiando $F_1$ y $F_2$ (y el signo con el que entran al tanque).

\paragraph{Escenario 3: prealimentación}
Si se puede medir $F_1$, se puede mover $F_2$ en la misma cantidad \emph{antes} de que el nivel cambie (\cref{fig:in-esquemas}(b)). Es la ventaja de la prealimentación: corrige antes del daño. Su debilidad es que no mira el nivel: cualquier error en la medición de $F_1$ o en el modelo se acumula sin que nadie lo corrija. Por eso en la práctica \emph{se combinan} las dos estructuras.

\begin{figure}[ht]
\centering
\includegraphics[width=0.92\linewidth]{cap01/esquema_escenarios}
\caption{Los escenarios 1 y 3 de la \cref{fig:in-escenarios} como diagramas de bloques.}
\label{fig:in-esquemas}
\end{figure}

\begin{ejemplo}[label=ej:in-prealimentacion]{Prealimentación con un sensor imperfecto}
Un tanque de área $A=2$\,m$^2$ cumple $A\,\dot h=F_1-F_2$ (la salida la impone una bomba). En $t=2$\,min, $F_1$ sube $0.5$\,m$^3$/min. Compare, en desviación $\dev{h}$ respecto al nivel inicial: (i) prealimentación con un transmisor de flujo que lee un $10\%$ menos, $\dev{F}_2=0.9\,\dev{F}_1$; (ii) retroalimentación proporcional, $\dev{F}_2=K_c\dev{h}$ con $K_c=1$\,m$^2$/min; (iii) las dos juntas.
\solucion (i) $A\dot{\dev{h}}=0.5-0.45=0.05$: el nivel sube $0.025$\,m/min \emph{para siempre}; en $18$\,min ya subió $0.45$\,m. (ii) $A\dot{\dev{h}}=0.5-K_c\dev{h}$: el nivel se estabiliza en $0.5/K_c=0.5$\,m con constante de tiempo $A/K_c=2$\,min. (iii) $A\dot{\dev{h}}=0.05-K_c\dev{h}$: se estabiliza en $0.05$\,m, diez veces menos que con retroalimentación sola. La \cref{fig:in-prealimentacion} lo muestra: la prealimentación hace casi todo el trabajo y la retroalimentación corrige lo que aquella no sabía.
\end{ejemplo}

\begin{figure}[ht]
\centering
\includegraphics[width=0.85\linewidth]{cap01/f_prealimentacion}
\caption{Ejemplo~\ref{ej:in-prealimentacion}: la prealimentación sola deriva, la retroalimentación sola deja una desviación grande, juntas funcionan.}
\label{fig:in-prealimentacion}
\end{figure}

\codigorepo{cap01_introduccion/C1_Ej3_Prealimentacion.m}{combinada}{ejemplo \ref{ej:in-prealimentacion}}

\begin{pausa}
En el escenario~1, si $F_2$ es una perturbación que \emph{se puede medir}, ¿cómo combinaría prealimentación y retroalimentación?
\respuesta{El controlador LC sigue corrigiendo con el nivel, y a su salida se le suma un término que mueve $F_1$ en la misma cantidad que cambie el $F_2$ medido. Así un cambio de $F_2$ se compensa de inmediato y LC solo corrige lo que la compensación no logró.}
\end{pausa}

%% ---------------------------------------------------------------------
\section{Una breve historia del control}
\label{sec:in-historia}
%% ---------------------------------------------------------------------

La historia del control automático ha sido contada con cuidado por Otto Mayr, para los orígenes~\citep{May71a,May12}; por Stuart Bennett, para el periodo industrial~\citep{Ben86,Ben93}; y por David Mindell, para los años de la guerra~\citep{Min02}. Hay números especiales de revistas dedicados al tema~\citep{Bus96,BeBu02}, una reseña breve y muy completa de Chris Bissell~\citep{Bis09} y una mirada de conjunto de Karl J.\ Åström y P.\,R.\ Kumar, \emph{Control: a perspective}~\citep{AsKu14}. Lo que sigue es un resumen de esas fuentes, con la referencia del trabajo original donde apareció cada idea cuando la hay. La \cref{fig:in-linea} resume la cronología.

\begin{figure}[ht]
\begin{anchocompleto}
\centering
\includegraphics[width=0.92\linewidth]{cap01/f_linea_tiempo}
\end{anchocompleto}
\caption{Algunos hitos de la historia del control. Arriba, de los reguladores de flotador a la teoría de reguladores y los pilotos automáticos; abajo, del amplificador de Black al control digital.}
\label{fig:in-linea}
\end{figure}

\notamargen{\textbf{Un patrón que se repite.} Primero llega el invento; luego aparece un problema de estabilidad que nadie entiende; y solo entonces llega la teoría. Watt $\to$ Maxwell, los barcos $\to$ Minorsky, el amplificador de Black $\to$ Nyquist.}

%% ---------------------------------------------------------------------
\subsection{De los relojes de agua a la máquina de vapor}

Como tecnología, el control tiene más de dos mil años. El primer mecanismo retroalimentado del que hay registro es el regulador de flotador de la \emph{clepsidra} (``ladrona de agua'') de Ktesibios de Alejandría, en el siglo~III a.C.~\citep{May71a}. Las sociedades que crecían necesitaban horarios, y los relojes de sol no sirven de noche ni con nubes. Un reloj de agua mide el tiempo por el volumen que cae a un recipiente, y para que sea preciso el caudal debe ser constante, lo que exige un nivel constante en el recipiente que gotea. Ktesibios lo logró con un flotador cónico que tapa la entrada cuando el nivel sube y la destapa cuando baja (\cref{fig:in-clepsidra}): sensor, comparador y actuador en una sola pieza.\nota{Ktesibios era hijo de un barbero y, según la tradición, empezó inventando un espejo ajustable para la barbería de su padre. Algunas fuentes lo hacen primer director del Museo de Alejandría.} Filón de Bizancio usó la misma idea, hacia el 250~a.C., para mantener constante el nivel de aceite de una lámpara.

\begin{figure}[ht]
\centering
\includegraphics[width=0.6\linewidth]{cap01/regulador_flotador}
\caption{El principio de la clepsidra de Ktesibios: el flotador mantiene constante el nivel del recipiente regulador, el caudal de salida es constante y el nivel del recipiente receptor sube de manera uniforme con el tiempo.}
\label{fig:in-clepsidra}
\end{figure}

En el siglo~I, Herón de Alejandría describió en su \emph{Pneumatica} varios reguladores de nivel y válvulas de flotador. Los ingenieros del mundo islámico los perfeccionaron: los hermanos Banū Mūsā en el siglo~IX y al-Jazarī en su \emph{Libro del conocimiento de los ingeniosos mecanismos} (1206) describen relojes y fuentes con reguladores de flotador~\citep{May71a}. En Europa, en cambio, estas ideas casi desaparecen entre la Edad Media y el Barroco: la \emph{Pneumatica} de Herón no se imprimió en latín sino hasta 1575~\citep{May12}.

El primer sistema retroalimentado de la Europa moderna que no copia los antiguos es el regulador de temperatura de Cornelis Drebbel (1572--1633), hacia 1620, para un horno que incubaba huevos: la dilatación del líquido de un termómetro empujaba una varilla que cerraba la entrada de aire al fuego~\citep{Bis09}. En 1681 Denis Papin agregó a su olla a presión una válvula de seguridad, antecedente de los reguladores de presión~\citep{May71a}.

\paragraph{Los molinos y el regulador centrífugo}
El siglo~XVIII trajo dos problemas de regulación en los molinos de viento: orientar las aspas hacia el viento y mantener la velocidad de las piedras. Para lo primero, Edmund Lee patentó en 1745 la \emph{cola de abanico}, un rotor auxiliar que gira el molino hasta enfrentarlo al viento. Para lo segundo se usaron péndulos centrífugos; Thomas Mead los patentó en 1787 para separar las piedras del molino según la velocidad~\citep{Bis09}. Hacia 1788 James Watt, a sugerencia de su socio Matthew Boulton, adaptó el regulador centrífugo a la máquina de vapor (\cref{fig:in-watt}): dos bolas que giran con el eje se abren cuando la máquina se acelera y, mediante un juego de palancas, cierran la válvula del vapor. El regulador combina sensado, actuación y control en un mismo mecanismo, y se convirtió en el símbolo de la Revolución Industrial. Watt no lo patentó: el principio ya se usaba en los molinos~\citep{May71a,Bis09}.

\begin{figure}[ht]
\centering
\includegraphics[width=\linewidth]{cap01/foto_watt}
\caption{Arriba a la izquierda, un modelo de la máquina de vapor de Watt; a la derecha, las bolas del regulador centrífugo; abajo, el autor junto a la máquina de vapor que conserva el Science Museum de Londres. Fotos del autor.}
\label{fig:in-watt}
\end{figure}

Los reguladores de Watt tenían un defecto: la velocidad de equilibrio cambiaba con la carga, porque la corrección era proporcional al exceso de velocidad. En el siglo~XIX se buscaron reguladores que llevaran la velocidad al mismo valor con cualquier carga. El \emph{regulador cronométrico} de los hermanos Werner y William Siemens lo logró con un engranaje diferencial que integraba la diferencia entre la velocidad de la máquina y la de un reloj: es la \clave{acción integral}~\citep{Bis09}. Además, a medida que las máquinas se volvieron más rápidas y los reguladores más sensibles, apareció un fenómeno inesperado, la \emph{caza} (\emph{hunting}): la máquina oscilaba alrededor de la velocidad deseada sin detenerse. George Biddell Airy, astrónomo real, publicó en 1840 el primer análisis matemático de un sistema retroalimentado al estudiar las oscilaciones del mecanismo que movía un telescopio~\citep{Air40}.

%% ---------------------------------------------------------------------
\subsection{La teoría de los reguladores (1868--1895)}

La teoría empieza con el artículo de James Clerk Maxwell \emph{On governors}, recibido por la Royal Society el 20 de febrero de 1868~\citep{Max68}. Maxwell escribió las ecuaciones diferenciales del regulador y de la máquina, las \emph{linealizó} alrededor del punto de operación y mostró que la estabilidad depende de que las raíces de la ecuación característica tengan parte real negativa. Dedujo las condiciones para sistemas de segundo y tercer orden, mostró las ventajas de la acción integral y distinguió entre los \emph{moderadores}, como el de Watt, que dejan un error que depende de la carga, y los verdaderos \emph{reguladores}, que no lo dejan~\citep{Kan16}. Es, en esencia, la diferencia entre la acción proporcional y la integral que veremos en los capítulos~6 y~7.

\begin{historia}[Maxwell, los anillos de Saturno y la estabilidad]
Antes de los reguladores, Maxwell había pensado en estabilidad por otra razón. En 1857 ganó el premio Adams de la Universidad de Cambridge con un ensayo sobre los anillos de Saturno, en el que descartó los modelos de anillos sólidos o fluidos porque sus movimientos no eran estables: si los anillos existen desde hace tanto tiempo, el modelo correcto tiene que ser estable. Usar la estabilidad como criterio para aceptar o rechazar un modelo es una idea que hoy usamos en identificación de sistemas~\citep{Bit15}. Norbert Wiener lo consideró el precursor de la teoría de la retroalimentación~\citep{May71b}.
\end{historia}

El artículo de Maxwell era difícil y tuvo poco efecto sobre los ingenieros de su tiempo; pero planteó una pregunta matemática: ¿cómo saber si todas las raíces de un polinomio de grado $n$ tienen parte real negativa sin calcularlas? Edward J.\ Routh la respondió y ganó con ello el premio Adams de 1877~\citep{Rou77}.\nota{Routh y Maxwell nacieron el mismo año, 1831, y presentaron juntos el examen de matemáticas de Cambridge (los \emph{tripos}) en 1854. Routh quedó primero (\emph{Senior Wrangler}) y Maxwell segundo. Ambos compartieron ese año el premio Smith, la primera vez que se dividió entre dos candidatos. (La anécdota del sirviente que, enviado a ver quién había quedado segundo, responde ``¡usted, señor!'' se cuenta de William Thomson ---lord Kelvin--- en 1845, no de Maxwell.)} En Rusia, Ivan Vyshnegradskii analizó en 1876 el regulador de la máquina de vapor sin conocer el trabajo de Maxwell, despreciando la fricción, linealizando y resolviendo el caso de tercer orden con un método gráfico~\citep{Vys76}; terminaría su carrera como ministro de finanzas del Imperio ruso. En Suiza, Aurel Stodola necesitaba en 1893 diseñar el regulador de una turbina hidráulica de alta presión para Davos; sus modelos eran de orden alto y le pidió ayuda a su colega del Politécnico de Zúrich, Adolf Hurwitz, quien publicó en 1895 un criterio equivalente al de Routh, escrito con determinantes~\citep{Bis89,Hur95}.\nota{Charles Hermite había resuelto en 1856 un problema matemático equivalente, sin relacionarlo con los reguladores.} En 1892 Aleksandr Lyapunov publicó su teoría de la estabilidad del movimiento, válida también para sistemas no lineales~\citep{Lya92}. Volveremos a todos ellos en el capítulo~5.

%% ---------------------------------------------------------------------
\subsection{Barcos, aviones, fábricas y teléfonos (1900--1940)}

El término \emph{servomoteur} lo acuñó en 1873 el ingeniero francés Joseph Farcot, que construía mecanismos hidráulicos para mover el timón de los barcos~\citep{Bis09}. A comienzos del siglo~XX el giroscopio cambió la navegación: Elmer Sperry fundó en 1910 la Sperry Gyroscope Company, que fabricó girocompases y luego pilotos automáticos para barcos que combinaban el giroscopio con un motor eléctrico conectado al timón. Su hijo Lawrence llevó la idea a los aviones: el 18 de junio de 1914, en el \emph{Concours de la Sécurité en Aéroplane} de París, voló con las manos en alto mientras su mecánico caminaba sobre el ala y el estabilizador giroscópico mantenía el avión nivelado~\citep{Bis09}.

Para gobernar barcos, Nicolas Minorsky observó a los timoneles de la Marina de Estados Unidos y en 1922 propuso una ley de control con tres términos: uno proporcional al error de rumbo, uno a su integral y uno a su derivada; estudió además la estabilidad del lazo con un modelo lineal~\citep{Min22,Ben84}. Es la primera formulación teórica del controlador \clave{PID}.

En la industria de procesos, los controladores neumáticos se volvieron comunes en los años veinte. La Foxboro anunció en 1931 el \emph{Stabilog}, cuya salida era proporcional al error y a su integral, y la Taylor Instrument incorporó en 1935 la acción derivativa (\emph{pre-act}) a su controlador Fulscope. Una encuesta del gobierno de Estados Unidos estimó que entre 1925 y 1935 se vendieron allí más de $75\,000$ controladores automáticos, y a mediados de los treinta había unas $600$ empresas fabricantes de instrumentos~\citep{BenPID93}. Faltaba saber cómo ajustar los parámetros: en 1942 John Ziegler y Nathaniel Nichols, de la Taylor, publicaron las reglas de sintonización que estudiaremos en el capítulo~8~\citep{ZiNi42}.

En los teléfonos estaba pasando algo parecido. Las líneas transcontinentales necesitaban cientos de amplificadores en cascada, cuya distorsión se acumulaba. Harold S.\ Black resolvió el problema en 1927 con el amplificador con retroalimentación negativa~\citep{black34}; la patente, presentada en 1928, tardó nueve años en concederse porque la idea parecía contraria a la intuición~\citep{blackpat,Bri11}. Con mucha ganancia de lazo los amplificadores oscilaban, y Harry Nyquist dio en 1932 el criterio para predecir la estabilidad del lazo cerrado a partir de la respuesta en frecuencia de la malla abierta~\citep{Nyq32,Bri10}. Hendrik Bode convirtió esas ideas, entre 1940 y 1945, en un método de diseño con márgenes de ganancia y de fase~\citep{Bod40,bode45}. El capítulo~\ref{cap:rf} cuenta esta parte en detalle.

En el MIT, mientras tanto, el grupo de Vannevar Bush construyó entre 1928 y 1931 el \emph{analizador diferencial}, un computador analógico mecánico para resolver ecuaciones diferenciales (\cref{fig:in-bush}). Harold Hazen, discípulo y colaborador de Bush durante más de veinte años, publicó en 1934 la \emph{Theory of servo-mechanisms}, la primera presentación sistemática de cómo funcionan y cómo se diseñan los servomecanismos~\citep{Haz34}. Su estudiante Gordon Brown dirigió desde los primeros años de la guerra el Laboratorio de Servomecanismos del MIT. La \cref{fig:in-personajes} reúne a varios de los protagonistas de esta sección.

\begin{figure}[ht]
\centering
\includegraphics[width=0.85\linewidth]{cap01/foto_bush}
\caption{Analizador diferencial del Science Museum de Londres, construido en Mánchester en 1935 siguiendo el diseño de Vannevar Bush en el MIT. Foto del autor.}
\label{fig:in-bush}
\end{figure}

\begin{figure}[ht]
\centering
\includegraphics[width=0.85\linewidth]{cap01/personajes}
\caption{Arriba, de izquierda a derecha: Edward J.\ Routh, Vannevar Bush, James C.\ Maxwell y Elmer A.\ Sperry. Abajo: Aurel Stodola, Nicolas Minorsky, Hendrik W.\ Bode y Adolf Hurwitz. Imágenes de Wikimedia Commons.}
\label{fig:in-personajes}
\end{figure}

%% ---------------------------------------------------------------------
\subsection{La guerra y la posguerra (1940--1960)}

La Segunda Guerra Mundial juntó a las tres comunidades: los ingenieros de procesos, los de servomecanismos y los de comunicaciones. Dirigir un cañón antiaéreo con un radar es un problema de servomecanismos con ruido, y en el Laboratorio de Radiación del MIT y en Bell Labs se unificaron los métodos de frecuencia. El resultado se publicó al terminar la guerra en el volumen \emph{Theory of Servomechanisms}~\citep{JaNiPh47}, el primer libro de texto moderno del área. En 1948 Norbert Wiener publicó \emph{Cybernetics}, que extendió la idea de retroalimentación a los seres vivos y a la sociedad~\citep{Wie48}, y Walter Evans presentó el lugar de las raíces~\citep{Eva50}.

Los computadores digitales traían otra pregunta: qué pasa cuando el controlador solo mide en ciertos instantes. En la Universidad de Columbia, John Ragazzini y Lotfi Zadeh introdujeron en 1952 la transformada $z$ para los sistemas muestreados~\citep{RaZa52}; por ese grupo pasaron también Eliahu Jury, Gene Franklin y Rudolf Kalman. La optimización avanzaba en paralelo: George Dantzig creó en 1947 el método símplex de programación lineal, y en la RAND Corporation Richard Bellman desarrolló en los años cincuenta la programación dinámica~\citep{Bel57}. En 1955 la Glenn L.\ Martin Company fundó en Baltimore el Research Institute for Advanced Studies (RIAS), donde Solomon Lefschetz reunió un grupo de matemáticos dedicados a la teoría de sistemas dinámicos y de control, entre ellos Kalman.\nota{Lefschetz tiene un lazo con la Universidad de los Andes. Mario Laserna, su fundador, organizó a finales de los años cuarenta en Princeton una reunión en honor de la nueva universidad, a la que asistieron Albert Einstein, John von Neumann, Oskar Morgenstern y Lefschetz, entre otros~\citep{IASLaserna}; en el verano de 1950 invitó a von Neumann y a Lefschetz a dictar conferencias en Bogotá~\citep{Les19}.}

El área se organizó institucionalmente. En 1945 se fundó en Estados Unidos la Instrument Society of America (ISA). En 1954 el Institute of Radio Engineers creó su grupo profesional de control automático, origen de la actual IEEE Control Systems Society. Y el 12 de septiembre de 1957, delegados de 18 países fundaron en París la International Federation of Automatic Control (IFAC); su primer presidente fue Harold Chestnut, y su primer congreso mundial se reunió en Moscú en 1960, en plena Guerra Fría. Desde entonces se reúne cada tres años~\citep{IFAChist,Bis09}.

%% ---------------------------------------------------------------------
\subsection{La edad de oro (1960--2000)}

La carrera espacial necesitaba guiar cohetes y satélites con muchas variables, restricciones de combustible y mediciones ruidosas, y los computadores empezaban a permitir cálculos que antes eran impensables. En el congreso de Moscú, Kalman presentó la teoría de sistemas en espacio de estado, con los conceptos de controlabilidad y observabilidad~\citep{Kal60}; el mismo año publicó el filtro que lleva su nombre~\citep{Kal60b}, extendido en 1961 con Richard Bucy~\citep{KaBu61}. El grupo de Lev Pontryagin publicó en 1962 el principio del máximo~\citep{Pon62}, que junto con la programación dinámica de Bellman fundó el \emph{control óptimo}. A esta manera de trabajar, con ecuaciones de estado en el dominio del tiempo, se le llamó \emph{control moderno}, para distinguirla del \emph{clásico} de Nyquist, Bode y Evans. El capítulo~\ref{cap:re} es nuestra puerta de entrada.

Del tronco común salieron muchas ramas: identificación de sistemas, control adaptativo, control no lineal, control robusto, sistemas de eventos discretos, control predictivo. En la industria el cambio fue el paso a lo digital. En 1968--1969 Modicon diseñó para la división Hydramatic de General Motors el primer controlador lógico programable (PLC), que reemplazó gabinetes enteros de relés; en 1975 Honeywell (TDC~2000) y Yokogawa (CENTUM) lanzaron los primeros sistemas de control distribuido (DCS); y la señal de $4$ a $20$\,mA reemplazó a la presión de $3$ a $15$\,psi de la instrumentación neumática. En los automóviles llegaron los frenos antibloqueo (ABS) y luego el control de estabilidad. El proyecto PATH de California, creado en 1986, demostró el 7 de agosto de 1997, en la autopista I-15 de San Diego, un pelotón de ocho automóviles que viajaban automáticamente a velocidad de autopista separados por unos $6$\,m~\citep{PATH97}.

%% ---------------------------------------------------------------------
\subsection{El control hoy}

Hoy el control está en casi todo, la mayoría de las veces escondido: en los discos duros y los teléfonos, en las redes eléctricas con fuentes renovables, en los procesos biológicos, en la logística y en los vehículos autónomos~\citep{AsKu14}. Los problemas nuevos tienen que ver con sistemas en red, con muchos agentes que deciden por su cuenta y con datos abundantes; pero las preguntas de fondo son las de Maxwell: ¿es estable?, ¿qué tan bien sigue lo que se le pide?, ¿qué tan robusto es ante lo que no sabemos? Quijano y Cotrino recogen la historia del control y la automatización en Colombia y el estado actual del área en el país~\citep{QuCo23}.

\begin{pausa}
¿Qué tienen en común la clepsidra de Ktesibios, el regulador de Watt y el amplificador de Black? Identifique en cada uno el sensor, el comparador y el actuador.
\respuesta{Los tres son lazos cerrados con retroalimentación negativa. Clepsidra: el flotador mide el nivel, su posición respecto al orificio hace de comparador y él mismo abre o cierra la entrada (actuador). Watt: las bolas miden la velocidad, el punto de equilibrio de las palancas fija la referencia y la válvula de vapor es el actuador. Black: la red resistiva $\beta$ mide la salida, la resta a la entrada se hace en el propio amplificador, y los tubos de vacío son el actuador.}
\end{pausa}

%% ---------------------------------------------------------------------
\section{El mapa del curso}
\label{sec:in-ruta}
%% ---------------------------------------------------------------------

Estas notas reúnen en un solo lugar lo que está disperso en muchos libros, con una visión propia, en castellano, que pone el énfasis en el análisis y en la simulación. No pretenden reemplazar a los textos clásicos: en cada capítulo se citan los libros y artículos de donde salen las ideas, para que usted pueda ir más lejos.

\subsection{La ruta}

El curso sigue el orden natural de un proyecto de control: \emph{modelar}, \emph{analizar}, \emph{diseñar} (\cref{fig:in-mapa}).

\begin{figure}[ht]
\centering
\includegraphics[width=\linewidth]{cap01/mapa_curso}
\caption{Los capítulos del curso y cómo dependen unos de otros. En azul, \emph{modelar}: la planta como ecuación diferencial y función de transferencia; en verde, \emph{analizar}: estabilidad, transitorio y error; en coral y violeta, \emph{diseñar}: el controlador.}
\label{fig:in-mapa}
\end{figure}

\begin{itemize}
  \item \textbf{Modelar} (capítulos~2 a~4). Para controlar un proceso hay que entenderlo, y lo entendemos con un modelo físico. Trabajaremos con sistemas \emph{lineales e invariantes en el tiempo} (LTI), descritos por ecuaciones diferenciales ordinarias. Los sistemas reales tienen parámetros que cambian, no linealidades y mediciones tomadas en instantes discretos; pero alrededor de un punto de operación, casi siempre un equilibrio, un modelo lineal describe bien su comportamiento. La transformada de Laplace convierte esas ecuaciones en funciones de transferencia, que se manipulan con álgebra, y la representación de estado las organiza para trabajar en el tiempo.
  \item \textbf{Analizar} (capítulos~5, 6 y~10). Con el modelo preguntamos si el lazo es estable (polinomio característico, valores propios, criterio de Routh--Hurwitz), cuánto error deja en estado estacionario y cómo responde a señales de distintas frecuencias (diagramas de Bode, márgenes de estabilidad, robustez).
  \item \textbf{Diseñar} (capítulos~7 a~10). El controlador PID, el más usado en la industria: qué hace cada término, sus configuraciones, sus problemas típicos (como la saturación del integrador) y cómo se sintonizan sus parámetros. La retroalimentación de estado, que ubica todos los polos y abre la puerta al control moderno, con los conceptos de controlabilidad y observabilidad. Y el diseño de compensadores en frecuencia con los diagramas de Bode.
\end{itemize}

\begin{receta}[Lo que usted sabrá hacer al final del curso]
\begin{enumerate}
  \item \textbf{Respuesta transitoria:} modelar, simular y analizar señales y sistemas en el tiempo (espacio de estado) y en la frecuencia (transformada de Laplace, función de transferencia).
  \item \textbf{Error en estado estacionario:} calcularlo según el tipo de entrada y el tipo del sistema, con las constantes de error $\kp$ y $\kv$.
  \item \textbf{Estabilidad:} decidirla con el polinomio característico o los valores propios, y estudiar cómo cambia con los parámetros (robustez).
  \item \textbf{Controladores:} diseñar PID y retroalimentaciones de estado, simularlos y mejorarlos (\emph{anti-windup}, acción integral, control por modelo interno, compensadores en frecuencia).
\end{enumerate}
\end{receta}

\subsection{Cómo usar estas notas}

Cada capítulo empieza con sus objetivos y un código QR que lleva a su página web, donde están el PDF, los códigos, los videos y un explorador interactivo. A lo largo del texto encontrará cajas de colores:

\begin{itemize}
  \item \textbf{Definición} y \textbf{Ejemplo} (numeradas): lo que hay que saber y cómo se usa. Los ejemplos tienen solución.
  \item \textbf{Idea clave}: el resumen de una sección en pocas líneas. \textbf{Cuidado}: un error frecuente.
  \item \textbf{Pausa activa}: una pregunta corta para comprobar que entendió; las respuestas están al final del capítulo.
  \item \textbf{Nota histórica}, \textbf{Para profundizar} y \textbf{Procedimiento}: contexto, material opcional y recetas paso a paso.
  \item \textbf{Código} y \textbf{Simulink}: fragmentos de los programas del repositorio, con un enlace al archivo completo. Antes de usarlos, ejecute \cmd{iniciarCurso} en \matlab. Cada carpeta tiene un \cmd{verificarCapNN} que comprueba los números del capítulo.
\end{itemize}

Los problemas del final de cada capítulo están marcados con \etiquetaM{} si requieren \matlab/\simulink y con \etiquetaEstrella{} si son de profundización. La notación es la misma en todo el curso y está resumida en la sección de convenciones al comienzo de las notas.

\subsection{Unas lecciones para volver a leer}

Dennis Bernstein escribió una guía de lecciones que todo estudiante de control debería ir afianzando~\citep{Ber97}. Muchas de sus palabras no tendrán sentido todavía; la idea es que vuelva a esta lista al terminar cada capítulo y vea cuántas ya entiende. Esta es una selección, adaptada al curso:

\begin{enumerate}
  \item \textbf{La retroalimentación está en todas partes}: en los procesos industriales, en la biología, en la economía.
  \item \textbf{Un diagrama de bloques no es un circuito.} Sus flechas son flujo de información, no conexiones físicas.
  \item \textbf{Encuentre los equilibrios y linealice.} Los sistemas reales son no lineales, pero el modelo lineal describe bien lo que pasa cerca del punto de operación.
  \item \textbf{Verifique primero la estabilidad}, con el polinomio característico o con los valores propios del modelo linealizado.
  \item \textbf{Los sistemas estables tienen respuesta en frecuencia}: una senoide de entrada produce, en régimen permanente, una senoide de la misma frecuencia con otra amplitud y otra fase.
  \item \textbf{No olvide la malla abierta} $\Gma(s)$: todo lo que hay en el lazo (controlador, actuador, planta, sensor, transmisor), su estabilidad y el signo de la retroalimentación.
  \item \textbf{Después de la estabilidad, el desempeño lo es todo.} Si la planta ya es estable, el desempeño es la única razón para retroalimentar.
  \item \textbf{La ganancia alta tiene beneficios}: rechaza perturbaciones y reduce la sensibilidad a los errores del modelo, como vimos en el tanque.
  \item \textbf{Pero tiene precios}: amplifica el ruido del sensor, satura los actuadores y puede desestabilizar el lazo.
  \item \textbf{Cuidado con los polos poco amortiguados}, que producen resonancias y oscilaciones largas.
  \item \textbf{Asegure la estabilidad nominal y luego la robusta}: que el lazo siga estable cuando la planta real no es la del modelo.
  \item \textbf{Respete lo inestable}~\citep{Ste03}: estabilizar una planta inestable impone límites a lo que se puede lograr.
  \item \textbf{Las no linealidades siempre están ahí}, el control multivariable no es trivial y, sobre todo, \textbf{la seguridad de las personas siempre está en juego}.
\end{enumerate}

%% ---------------------------------------------------------------------
\section{Para llevar}
\label{sec:in-cierre}
%% ---------------------------------------------------------------------

\begin{parallevar}
\begin{enumerate}
  \item Un sistema de control administra el comportamiento de otro sistema para que una salida siga una referencia o rechace perturbaciones. Analizar es determinar su desempeño; diseñar es lograrlo.
  \item Retroalimentar es medir, comparar con lo deseado y actuar sobre el error. Con una ganancia de lazo $KK_c$, las perturbaciones y los errores del modelo se reducen aproximadamente por $1+KK_c$. El precio: esfuerzo de control, ruido de medición y riesgo de inestabilidad.
  \item El control en lazo abierto es exacto solo si el modelo es exacto y no hay perturbaciones. La prealimentación corrige antes del daño; la retroalimentación corrige lo que no se sabía. En la práctica se combinan.
  \item ``Control en lazo abierto'' es una estrategia; la \emph{malla abierta} $\Gma=G\Hs$ y el \emph{lazo abierto} $\Gla$ son funciones de transferencia de un lazo cerrado.
  \item El desempeño se mide con la estabilidad, la respuesta transitoria y el error en estado estacionario.
  \item Un problema de control se plantea con objetivos, variables manipuladas y perturbaciones, salidas medidas y no medidas, restricciones, tipo de operación, consideraciones de seguridad y estructura del controlador.
  \item La teoría del control nació casi siempre de un problema de estabilidad: el regulador de Watt llevó a Maxwell, Routh y Hurwitz; los barcos, a Minorsky; los amplificadores, a Nyquist y Bode; la carrera espacial, a Kalman.
\end{enumerate}
\end{parallevar}

\begin{profundiza}[Correcciones respecto a la versión anterior del capítulo]
(i) Watt no patentó el regulador centrífugo; Mead lo había patentado para los molinos en 1787. (ii) El vuelo de Lawrence Sperry en París fue en 1914, no en 1912. (iii) Minorsky propuso el PID y analizó su estabilidad con un modelo lineal; no usó ubicación de polos. (iv) Bode publicó sus resultados entre 1940 y 1945, no en 1934; ese año fue el de la teoría de servomecanismos de Hazen (antes ``Hanzen''). (v) IFAC se fundó en 1957 con 18 países; RIAS en 1955. (vi) ``Stabilorg'' $\to$ Stabilog; ``Danzig'' $\to$ Dantzig; ``Vaneevar'' $\to$ Vannevar. (vii) Se retiró la atribución de la acción integral a los hermanos Périer, que no se encontró en las fuentes; se cita en su lugar el regulador cronométrico de los Siemens. (viii) El ejemplo del tanque ahora compara el lazo abierto con el cerrado y tiene números verificados con \texttt{herramientas/cap01\_numeros.py} y \cmd{verificarCap01.m}. (ix) La función $G(s)H(s)$ se llama \emph{malla abierta}, no ``función de transferencia en lazo abierto''.
\end{profundiza}

%% ---------------------------------------------------------------------
\section{Problemas}
\label{sec:in-problemas}
%% ---------------------------------------------------------------------

Los problemas marcados con \etiquetaM{} requieren \matlab/\simulink; los marcados con \etiquetaEstrella{} son de profundización.

\begin{problema}[lazo abierto o cerrado]\label{prob:in-clasificar}
Clasifique como control en lazo abierto o en lazo cerrado y, en los cerrados, identifique la variable controlada, el sensor, el actuador y una perturbación: (a) una tostadora con temporizador; (b) la nevera de su casa; (c) un semáforo de tiempos fijos; (d) un semáforo que cuenta los carros con una cámara; (e) la regulación de la temperatura del cuerpo humano; (f) una persona que se ducha (\cref{fig:in-ducha}); (g) el control de crucero de un carro. ¿Alguno usa prealimentación?
\end{problema}

\begin{problema}[el tanque, en detalle]\label{prob:in-tanque}
\etiquetaM{} Para el tanque del ejemplo~\ref{ej:in-tanque}: (a) reemplace la ley de lazo cerrado en \eqref{eq:in-tanque} y obtenga la ecuación diferencial del lazo cerrado y su constante de tiempo. (b) Calcule el nivel final con lluvia y con la válvula obstruida para $K_c=1$ y $K_c=10$, y compare con el factor $1+KK_c$. (c) Verifique (b) con \cmd{C1\_Ej1\_TanqueLazos}. (d) Limite el caudal a $0\le u\le2$\,m$^3$/min y repita la simulación con $K_c=10$: ¿qué cambia en el transitorio? ¿Y en el valor final? (e) Quite el término de calibración, $u=K_c(h_r-h)$: ¿cuánto vale el nivel final sin perturbación con $K_c=5$? ¿Por qué?
\end{problema}

\begin{problema}[plantear un problema]\label{prob:in-plantear}
Escoja uno: (a) un invernadero en la sabana de Bogotá; (b) un calentador de agua a gas; (c) la dosificación de insulina en una persona con diabetes. Plantee el problema de control con la receta de la \cref{sec:in-problema}: objetivos, manipuladas, perturbaciones, salidas medidas y no medidas, restricciones duras y blandas, tipo de operación, seguridad y una estructura de control con un diagrama de bloques.
\end{problema}

\begin{problema}[ganancia del proceso]\label{prob:in-ganancia}
(a) Un calentador eléctrico pasa de $2.0$ a $2.5$\,kW y la temperatura de un fluido pasa de $40$ a $50$\,°C: halle la ganancia con unidades. (b) En un enfriador, abrir la válvula de agua fría del $30$ al $40\%$ baja la temperatura de $25$ a $21$\,°C: halle la ganancia. ¿Qué significa su signo para el controlador? (c) Si la ganancia de (a) cambia a la mitad con la temperatura ambiente, ¿por qué eso complica el control en lazo abierto?
\end{problema}

\begin{problema}[prealimentación]\label{prob:in-prealimentacion}
\etiquetaM{} Para el ejemplo~\ref{ej:in-prealimentacion}: (a) resuelva analíticamente las tres ecuaciones diferenciales y verifique los valores de la \cref{fig:in-prealimentacion}. (b) ¿Qué error del transmisor de flujo se puede tolerar con prealimentación sola si la deriva debe ser menor que $1$\,cm/min? (c) Con las dos estructuras combinadas, ¿qué $K_c$ deja una desviación final menor que $1$\,cm? (d) Verifique con \cmd{C1\_Ej3\_Prealimentacion}.
\end{problema}

\begin{problema}[moderadores y reguladores]\label{prob:in-maxwell}
\etiquetaEstrella{} Maxwell distinguía los \emph{moderadores}, que dejan un error que depende de la carga, de los \emph{reguladores}, que no lo dejan. Para el tanque del ejemplo~\ref{ej:in-tanque} considere $u=h_r/K+K_ce+K_i\int_0^te\,\dif\tau$, con $e=h_r-h$. (a) Suponiendo que el lazo llega a un equilibrio, muestre que en él $h=h_r$ para cualquier lluvia constante y cualquier $R$. (b) ¿Cuál de las dos clases de Maxwell es el controlador del ejemplo~\ref{ej:in-tanque}? (c) \etiquetaM{} Simule con $K_c=5$ y $K_i=10$ y compare con la \cref{fig:in-tanque}.
\end{problema}

\begin{problema}[el ascensor]\label{prob:in-ascensor}
\etiquetaM{} Con \cmd{C1\_Ej2\_Ascensor}, cambie el factor de amortiguamiento a $\zeta=1$ y a $\zeta=0.3$ manteniendo $\wn=0.8$\,rad/s. Reporte el sobrepaso y el tiempo de establecimiento ($2\%$) y discuta cuál ascensor preferiría usted como pasajero y cuál como persona apurada.
\end{problema}

\begin{problema}[un invento, una fecha, una referencia]\label{prob:in-historia}
\etiquetaEstrella{} Escoja un dispositivo o una idea de la \cref{sec:in-historia}, busque la publicación o patente original y escriba una página que responda: ¿qué problema resolvía?, ¿dónde estaban el sensor, el comparador y el actuador?, ¿qué problema de estabilidad o de desempeño apareció?, ¿quién lo resolvió y cómo? Cite la fuente original y al menos una reseña histórica de las de esta sección.
\end{problema}

%% ---------------------------------------------------------------------
\section{Soluciones seleccionadas}
%% ---------------------------------------------------------------------

\paragraph{Problema~\ref{prob:in-clasificar}}
Lazo abierto: (a) y (c). Lazo cerrado: (b) temperatura, termostato, compresor, abrir la puerta; (d) longitud de las colas, cámara, tiempos del semáforo, llegada de carros; (e) temperatura corporal, termorreceptores, sudor y escalofríos, el clima; (f) y (g) velocidad, sensor de velocidad, acelerador, pendientes y viento. La ducha usa además prealimentación cuando uno oye el sanitario.

\paragraph{Problema~\ref{prob:in-tanque}}
(a) $\tau\dot h+(1+KK_c)h=K\bigl(h_r/K+K_ch_r+d\bigr)$; constante de tiempo $\tau/(1+KK_c)$. (b) Con $K_c=1$: lluvia $h=1.133$, obstrucción $h=1.059$; con $K_c=10$: $1.019$ y $1.008$. La desviación por lluvia es $0.4/(1+2K_c)$. (d) El caudal se queda en $2$ al comienzo y el nivel sube más despacio; el valor final no cambia, porque en equilibrio el caudal pedido es menor que $2$. (e) $h=RK_c/(1+RK_c)=10/11=0.909$: sin calibración el controlador proporcional necesita un error para producir caudal.

\paragraph{Problema~\ref{prob:in-ganancia}}
(a) $10/0.5=20$\,°C/kW. (b) $-4/10=-0.4$\,°C por punto porcentual: la ganancia es negativa, y el controlador debe abrir la válvula cuando la temperatura está \emph{por encima} de la referencia (acción inversa). (c) La calibración $u=r/K$ queda mal y el error es del $100\%$ del cambio.

\paragraph{Problema~\ref{prob:in-prealimentacion}}
(a) $\dev{h}=0.025(t-2)$; $\dev{h}=0.5\bigl(1-e^{-(t-2)/2}\bigr)$; $\dev{h}=0.05\bigl(1-e^{-(t-2)/2}\bigr)$, para $t\ge2$. (b) $0.5\,\varepsilon/2<0.01$: error menor que el $4\%$. (c) $0.05/K_c<0.01$: $K_c>5$\,m$^2$/min.

\paragraph{Problema~\ref{prob:in-maxwell}}
(a) En equilibrio todas las señales son constantes; si $\int e$ es constante, su derivada $e$ es cero, luego $h=h_r$, sin importar $d$ ni $R$ (siempre que el lazo sea estable). (b) Un moderador: deja un error que depende de la perturbación.

\paragraph{Problema~\ref{prob:in-ascensor}}
$\zeta=1$: sin sobrepaso, $\ts=7.29$\,s. $\zeta=0.6$: $9.48\%$ y $7.43$\,s. $\zeta=0.3$: $37.2\%$ y $14.0$\,s. Menos amortiguamiento no siempre es más rápido: con $\zeta=0.3$ llega antes al piso ($\tp=4.12$\,s, contra $4.91$\,s) pero tarda el doble en quedarse quieto.

\respuestaspausas

